\exercisesfor{3}

\begin{enumerate}
\def\labelenumi{\arabic{enumi}.}
\item
  Let \(G\) be a Lie group and \(G_1\) and \(G_2\) Lie subgroups of \(G\) . Suppose that \(G_2\) is of finite codimension in \(G\) and that \(K\) is of characteristic 0. Then \(G_1 \cap G_2\) is a Lie subgroup of \(G\) with Lie algebra \(L(G_1) \cap L(G_2)\) (argue as in Proposition 29 and its Corollary 2).
\item
  Suppose that \(\mathbf{K}\) is locally compact. Let \(\mathbf{G}\) be a Lie group of finite dimension \(n\) .
\end{enumerate}

\begin{enumerate}
\def\labelenumi{(\alph{enumi})}
\tightlist
\item
  Let \(\phi\) be an étale endomorphism of G with finite kernel. Let \(\mu\) be a left Haar measure of G. Then \(\phi\) is proper and
\end{enumerate}

\[
\phi (\mu) = \alpha \cdot \mu | \phi (G) \quad \text { where } \quad \alpha = \frac {\operatorname{Card} (\operatorname{Ker} \phi)}{\operatorname{mod} \det L (\phi)}.
\]

(Use Proposition 55.)
% semantic-fragments: legacy-input-seam
\relax{}
\begin{enumerate}
\def\labelenumi{(\alph{enumi})}
\setcounter{enumi}{1}
\tightlist
\item
  Suppose that G is compact. Let \(\phi\) be an étale endomorphism of G. Then \(\operatorname{Ker} \phi\) is finite, \(\phi(G)\) is of finite index in G and
\end{enumerate}

\[
\frac {\operatorname{Card} (\operatorname{Ker} \phi)}{\operatorname{Card} (\mathrm{G} / \phi (\mathrm{G}))} = \operatorname{mod} \det \mathrm{L} (\phi).
\]

(Use (a) to transform \(\mu (\mathbf{G})\) .)

\begin{enumerate}
\def\labelenumi{(\alph{enumi})}
\setcounter{enumi}{2}
\tightlist
\item
  Suppose that G is compact and commutative. Let \(r \in \mathbf{Z}\) be such that \(r, 1 \neq 0\) in K and let \(\phi\) be the endomorphism \(x \mapsto x^r\) of G. Then
\end{enumerate}

\[
\frac {\operatorname{Card} (\operatorname{Ker} \phi)}{\operatorname{Card} (\mathrm{G} / \phi (\mathrm{G}))} = (\mathrm{mod} r. 1) ^ {n}.
\]

(Observe that \(\mathrm{L}(\phi)\) is the homothety of ratio \(r\) by §2, Corollary to Proposition 7.)

\begin{enumerate}
\def\labelenumi{\arabic{enumi}.}
\setcounter{enumi}{2}
\tightlist
\item
  Let E be the complex vector space of symmetric complex matrices with 2 rows and 2 columns. For all \(s \in \mathbf{SL}(2, \mathbf{C})\) , we define an automorphism \(\varphi(s)\) of E by \(\varphi(s).m = s.m.^t s\) .
\end{enumerate}

\begin{enumerate}
\def\labelenumi{(\alph{enumi})}
\tightlist
\item
  Show that \(\rho\) is an analytic linear representation of \(\mathbf{SL}(2, \mathbf{C})\) . Identifying \(\begin{pmatrix} a & b \\ b & c \end{pmatrix} \in \mathrm{E}\) with \((a, b, c) \in \mathbf{C}^3\) , the matrix of \(\rho\begin{pmatrix} \alpha & \beta \\ \gamma & \delta \end{pmatrix}\) is
\end{enumerate}

\[
\left( \begin{array}{c c c} \alpha^ {2} & 2 \alpha \beta & \beta^ {2} \\ \alpha \gamma & \alpha \delta + \beta \gamma & \beta \delta \\ \gamma^ {2} & 2 \gamma \delta & \delta^ {2} \end{array} \right).
\]

\begin{enumerate}
\def\labelenumi{(\alph{enumi})}
\setcounter{enumi}{1}
\tightlist
\item
  Show that \(m \mapsto \det(m)\) is the non-degenerate quadratic form
\end{enumerate}

\[
q (a, b, c) = a c - b ^ {2}
\]

on E and that \(\rho\) is a morphism of \(\mathbf{SL}(2,\mathbf{C})\) onto \(\mathbf{SO}(q)\) with kernel \(\{I, - I\}\) . (For the surjectivity, observe that, in the matrix of an element of \(\mathbf{SO}(q)\) , the 1st column can be expressed in the form \((\alpha^2,\alpha \gamma ,\gamma^2)\) and the 3rd column in the form \((\beta^2,\beta \delta ,\delta^2)\) , where \(\alpha^2\delta^2 +\beta^2\gamma^2 -2\alpha \gamma \beta \delta = 1\) . Then observe that 2 elements of \(\mathbf{SO}(q)\) which coincide on a non-isotropic plane are equal.)

\begin{enumerate}
\def\labelenumi{(\alph{enumi})}
\setcounter{enumi}{2}
\tightlist
\item
  Deduce that the complex Lie group \(\mathbf{SO}(3, \mathbf{C})\) is isomorphic to the complex Lie group \(\mathbf{SL}(2, \mathbf{C}) / \{I, -I\}\) .
\end{enumerate}

\begin{enumerate}
\def\labelenumi{\arabic{enumi}.}
\setcounter{enumi}{3}
\tightlist
\item
  \begin{enumerate}
  \def\labelenumii{(\alph{enumii})}
  \tightlist
  \item
    Let \(q\) be a quadratic form of signature (1, 2) on \(\mathbf{R}^3\) . Let \(\mathbf{P}\) be the set of \(m \in \mathbf{R}^3\) such that \(q(m) > 0\) . Then \(\mathbf{P}\) is the union of two disjoint convex cones \(\mathbf{C}\) and \(-\mathbf{C}\) . If \(s \in \mathbf{SO}(q)\) , then either \(s(\mathbf{C}) = \mathbf{C}\) and \(s(-\mathbf{C}) = -\mathbf{C}\) or \(s(\mathbf{C}) = -\mathbf{C}\) and \(s(-\mathbf{C}) = \mathbf{C}\) . Let \(\mathbf{SO}^{+}(q)\) be the set of \(s \in \mathbf{SO}(q)\) such that \(s(\mathbf{C}) = \mathbf{C}\) . Then \(\mathbf{SO}^{+}(q)\) is an open normal subgroup of \(\mathbf{SO}(q)\) of index 2 in \(\mathbf{SO}(q)\) .
  \end{enumerate}
\end{enumerate}

\begin{enumerate}
\def\labelenumi{(\alph{enumi})}
\setcounter{enumi}{1}
\item
  By imitating the method of Exercise 3, define a morphism of \(\mathbf{SL}(2,\mathbf{R})\) onto \(\mathbf{SO}^{+}(q)\) with kernel \(\{I, - I\}\) . Deduce that the real Lie group \(\mathbf{SO}^{+}(q)\) is isomorphic to the real Lie group \(\mathbf{SL}(2,\mathbf{R}) / \{I, - I\}\) .
\item
  For \((\xi, \eta) \in \mathbf{C}^2\) and \((\xi', \eta') \in \mathbf{C}^2\) , let \(f((\xi, \eta), (\xi', \eta')) = \xi \bar{\xi}' - \eta \bar{\eta}'\) . Show that the inner automorphism of \(\mathbf{SL}(2, \mathbf{C})\) defined by \(\frac{1}{\sqrt{2}} \begin{pmatrix} 1 & -i \\ -i & 1 \end{pmatrix}\) maps \(\mathbf{SU}(f)\) to \(\mathbf{SL}(2, \mathbf{R})\) .
\end{enumerate}

\begin{enumerate}
\def\labelenumi{\arabic{enumi}.}
\setcounter{enumi}{4}
\tightlist
\item
  Let E be the real vector space of Hermitian complex matrices with 2 rows and 2 columns of zero trace. For all \(s \in \mathbf{SU}(2, \mathbf{C})\) , we define an automorphism \(\rho(s)\) of E by \(\rho(s)(m) = sms^*\) . We identify the element
\end{enumerate}

\[
\left( \begin{array}{c c} x & y + i z \\ y - i z & - x \end{array} \right)
\]

of E with the element \((x,y,z)\) of \(R^{3}\) . By imitating the method of Exercise 3, show that \(\rho\) is a morphism of the real Lie group \(\mathbf{SU}(2,\mathbf{C})\) onto the real Lie group \(\mathbf{SO}(3,\mathbf{R})\) with kernel \(\{I,-I\}\) and that \(\mathbf{SO}(3,\mathbf{R})\) is isomorphic to \(\mathbf{SU}(2,\mathbf{C})/\{I,-I\}\)

¶ 6. Let F be the vector space of Hermitian complex matrices with 2 rows and 2 columns. For all \(s \in \mathbf{SL}(2, \mathbf{C})\) , we define an automorphism \(\sigma(s)\) of F by \(\sigma(s)(m) = sms^*\) . We identify the element \(\begin{pmatrix} t + x & y + iz \\ y - iz & t - x \end{pmatrix}\) of F with the element \((t, x, y, z)\) of \(\mathbf{R}^4\) and write

\[
q (t, x, y, z) = t ^ {2} - x ^ {2} - y ^ {2} - z ^ {2}.
\]

Let C be the set of \((t, x, y, z) \in \mathbf{R}^+\) such that \(t^2 - x^2 - y^2 - z^2 > 0\) , \(t > 0\) . Let \(\mathbf{SO}^+(q)\) be the set of \(g \in \mathbf{SO}(q)\) such that \(g(\mathbf{C}) = \mathbf{C}\) ; it is an open normal subgroup of \(\mathbf{SO}(q)\) of index 2 in \(\mathbf{SO}(q)\) (cf.~Exercise 4). Show that \(\sigma\) is a morphism of the underlying real Lie group of \(\mathbf{SL}(2, \mathbf{C})\) onto the real Lie group \(\mathbf{SO}^+(q)\) with kernel \(\{I, -I\}\) . (For the surjectivity, use Exercise 5.) Deduce that the real Lie group \(\mathbf{SO}^+(q)\) is isomorphic to the underlying real Lie group of \(\mathbf{SL}(2, \mathbf{C}) / \{I, -I\}\) , that is (Exercise 3) the underlying real Lie group of \(\mathbf{SO}(3, \mathbf{C})\) .

\begin{enumerate}
\def\labelenumi{\arabic{enumi}.}
\setcounter{enumi}{6}
\tightlist
\item
  \begin{enumerate}
  \def\labelenumii{(\alph{enumii})}
  \tightlist
  \item
    Let G be the set of quaternions of norm 1. This is a Lie subgroup of the real Lie group \(\mathbf{H}^*\) , homeomorphic to \(\mathbf{S}_3\) and hence simply connected, cf.~General Topology, Chapter XI.
  \end{enumerate}
\end{enumerate}

\begin{enumerate}
\def\labelenumi{(\alph{enumi})}
\setcounter{enumi}{1}
\tightlist
\item
  Let E be the real vector space of pure quaternions, identified with \(\mathbf{R}^3\) by
\end{enumerate}

\[
(x, y, z) \mapsto x i + y j + z k.
\]

For all \(g \in \mathbf{G}\) , we define an automorphism \(\rho(g)\) of \(\mathbf{E}\) by \(\rho(g)q = ggg^{-1}\) . Show that \(\rho\) is a morphism of the Lie group \(\mathbf{G}\) onto the real Lie group \(\mathbf{SO}(3, \mathbf{R})\) with kernel \(\{1, -1\}\) . Deduce that the Lie group \(\mathbf{SO}(3, \mathbf{R})\) is isomorphic to the Lie group \(G / \{1, -1\}\) .

\begin{enumerate}
\def\labelenumi{(\alph{enumi})}
\setcounter{enumi}{2}
\item
  Let \(\mathbf{C}\) be identified with the subfield \(\mathbf{R} + \mathbf{R}i\) of \(\mathbf{H}\) . The mapping \((\lambda, q) \mapsto q\lambda\) of \(\mathbf{C} \times \mathbf{H}\) into \(\mathbf{H}\) makes \(\mathbf{H}\) into a vector space over \(\mathbf{C}\) which is identified with \(\mathbf{C}^2\) by the choice of basis \((1, j)\) . For all \(g \in \mathbf{G}\) , let \(\sigma(g)\) be the automorphism of this vector space defined by \(\sigma(g)q = gq\) . Show that \(\sigma\) is an isomorphism of the real Lie group G onto the real Lie group \(\mathbf{SU}(2, \mathbf{C})\) . The latter is therefore simply connected by \((a)\) .
\item
  For \((q_{1}, q_{2}) \in \mathbf{G} \times \mathbf{G}\) , let \(\tau(q_{1}, q_{2})\) be the mapping \(q \mapsto q_{1}q\bar{q}_{2}\) of \(\mathbf{H} = \mathbf{R}^{4}\) into itself. Show that \(\tau\) is a morphism of the real Lie group \(\mathrm{G} \times \mathrm{G}\) onto the real Lie group \(\mathbf{SO}(4, \mathbf{R})\) with kernel \(\mathrm{N} = \{(1, 1), (-1, -1)\}\) and hence that \(\mathbf{SO}(4, \mathbf{R})\) is isomorphic to \((\mathrm{G} \times \mathrm{G}) / \mathrm{N}\) .
\end{enumerate}

\begin{enumerate}
\def\labelenumi{\arabic{enumi}.}
\setcounter{enumi}{7}
\tightlist
\item
  Let \(\mathrm{G}\) be a finite-dimensional connected real Lie group and \(\mathbf{M}\) a finite-dimensional connected manifold of class \(\mathrm{C}^{\infty}\) ; let there be given a law of right operation of class \(\mathrm{C}^{\infty}\) of \(\mathrm{G}\) on \(\mathbf{M}\) . For all \(x \in \mathbf{M}\) , let \(\varphi(x)\) be the corresponding orbital mapping.
\end{enumerate}

\begin{enumerate}
\def\labelenumi{(\alph{enumi})}
\item
  For G to operate transitively on M, it is necessary and sufficient that, for all \(x \in M\) , the mapping \(\mathrm{T}_{e}(\varphi(x))\) of L(G) into \(\mathrm{T}_{x}(\mathrm{M})\) be surjective. (To see that the condition is sufficient, observe that then every G-orbit in M is open.)
\item
  Suppose that G operates transitively on M. We identify M with a homogeneous space H\textbar G, where H is the stabilizer of a point \(x_{0}\) of M. The action of G on M is called imprimitive if there exists a closed submanifold V of M of class \(C^{\omega}\) such that \(0 < dim V < dim M\) and for which every transform V.s (where \(s \in G\) ) is either equal to V or disjoint from V. For this to be so, it is necessary and sufficient that there exist a Lie subgroup L of G such that \(H \subset L \subset G\) and \(\dim(H) < \dim(L) < \dim(G)\) . If this is not so, the action of G on M is called primitive; for this it suffices that there exist no subalgebra of L(G) lying between L(G) and L(H) and distinct from L(G) and L(H).
\item
  Under the hypothesis of (b), let \(\mathcal{L}\) be the image of L(G) under the law of infinitesimal operation associated with the action of G. Let \(\mathcal{E}\) be the algebra of functions of class \(C^\infty\) on M and m the maximal ideal of \(\mathcal{E}\) consisting of the functions in \(\mathcal{E}\) which are zero at \(x_0\) ; we denote by \(\mathcal{L}_p\) for \(p = -1, 0, 1, 2, \ldots\) the set of \(x \in \mathcal{L}\) such that \(\theta(\mathbf{X}) f \in m^{p+1}\) for all \(f \in \mathcal{E}\) ; then \(\mathcal{L}_{-1} = \mathcal{L}\) and \(\mathcal{L}_0\) is the image of L(H) under the law of infinitesimal operation. If (U, \(\phi, \mathbf{R}^n\) ) is a chart of M centered on \(x_0\) and \(\mathbf{X}_1, \ldots, \mathbf{X}_n\) are the vector fields on U defined by the chart, the elements of \(\mathcal{L}_p\) are the \(a_1\mathbf{X}_1 + \cdots + a_n\mathbf{X}_n\) with \(a_1, \ldots, a_n\) in \(m^p|\text{U}\) . Show that \([\mathcal{L}_p, \mathcal{L}_q] \subset \mathcal{L}_{p,q}\) (with the convention \(\mathcal{L}_{-2} = \mathcal{L}\) ). If there exists Y \(\in \mathcal{L}_p\) (for \(p \geqslant 0\) ) such that Y \(\notin \mathcal{L}_{p+1}\) , there exists X \(\in\) L such that {[}Y, X{]} \(\notin \mathcal{L}_p\) . The \(\mathcal{L}_p\) are Lie subalgebras of \(\mathcal{L}\) . For \(p \geqslant 0\) , \(\mathcal{L}_p\) is an ideal of \(\mathcal{L}_0\) . Let \(\rho\) be the canonical linear representation of H on T \(_{x_0}\) (M). The Lie algebra of H/(Ker \(\rho\) ) is isomorphic to \(\mathcal{L}_0/\mathcal{L}_1\) .
\item
  Suppose further that the intersection of the \(\mathcal{L}_p\) is \(\{0\}\) . Then there exists a greatest index \(r\) such that \(\mathcal{L}_r \neq \{0\}\) and all the \(\mathcal{L}_f\) of indices \(\leqslant r\) are distinct. For \(p \geqslant 0\) ,
\end{enumerate}

\[
\dim (\mathcal {L} _ {p} / \mathcal {L} _ {p + 1}) \leqslant \binom {n + p} {p + 1}.
\]

\begin{enumerate}
\def\labelenumi{(\alph{enumi})}
\setcounter{enumi}{4}
\tightlist
\item
  The hypothesis of (d) is satisfied if M is analytic and G operates analytically on M.
\end{enumerate}

\begin{enumerate}
\def\labelenumi{\arabic{enumi}.}
\setcounter{enumi}{8}
\tightlist
\item
  In the notation of Proposition 30, HH' may be an open subset of G dense in G and distinct from G (take G = GL(2, R) with H the upper triangular subgroup and H' the lower strict triangular subgroup).
\end{enumerate}

